% !TeX root = ../main.tex

\chapter{Control Architecture}\label{chapter:control-architecture}

This chapter presents the flight control system used on the X-wing quadrotor.
The architecture combines (i) \emph{flatness-based reference generation} for dynamically
feasible trajectories, (ii) a \emph{geometric outer loop} on $SE(3)$ that converts
position and velocity errors into desired attitude and collective thrust, and
(iii) an \emph{Incremental Nonlinear Dynamic Inversion (INDI)} inner loop that
realizes the requested angular dynamics while rejecting unmodeled aerodynamic
effects from the wings.
The wings are not explicitly modeled; their influence is treated as a disturbance
and compensated incrementally~\cite{Sieberling2010,vanKampen2018,Tzoumanikas2021}.

\section{Reference Generation}\label{sec:ref-gen}
Quadrotor dynamics are differentially flat with respect to
$\mathbf{p}_d=[x_d,y_d,z_d]^\top$ and $\psi_d$.
We generate $(\mathbf{p}_d,\dot{\mathbf{p}}_d,\ddot{\mathbf{p}}_d,\psi_d)$
from a minimum-snap polynomial~\cite{Mellinger2011}.
In forward flight, sideslip is reduced through a coordinated-turn law,
\[
\dot{\psi}_d \approx \frac{a_{y,d}}{V_d\cos\theta_d},
\]
which aligns the body $x$-axis with the velocity vector.

\section{Geometric Outer Loop}\label{sec:outer}
The outer loop maps position and velocity errors to a desired attitude $R_d\in SO(3)$
and collective thrust $f_{\text{coll}}$~\cite{Goodarzi2015}.
Nominal commanded acceleration is
\begin{equation}
\mathbf{a}_c = \mathbf{k}_p\odot(\mathbf{p}_d-\mathbf{p})
              + \mathbf{k}_v\odot(\dot{\mathbf{p}}_d-\mathbf{v})
              + \ddot{\mathbf{p}}_d - \mathbf{g}.
\label{eq:ac-nom}
\end{equation}

\section{Incremental Aerodynamic Compensation (Outer INDI Layer)}\label{sec:outer-indi}
The controller estimates a residual between the thrust-induced specific force
(from filtered motor speeds) and the IMU-measured specific force.
That residual is rotated into the world frame and added to the commanded acceleration:
\begin{equation}
\mathbf{a}_{\text{aero}} = R\!\left(\frac{f_{T,f}}{m}\mathbf{e}_3 - \mathbf{a}_{B,f}\right),
\qquad
\mathbf{a}_c \leftarrow \mathbf{a}_c + \mathbf{a}_{\text{aero}}
\quad \text{if } z>0.5~\mathrm{m}.
\label{eq:aero-comp-outer}
\end{equation}
This constitutes an \emph{incremental} aerodynamic feedforward identical in principle
to INDI: rather than modeling lift or drag explicitly, the controller compensates
their measured effect from the previous time step.
As a result, quasi-steady aerodynamic forces vanish from the translational control law.

The desired body $z$-axis is set to $\mathbf{z}_B^d=\mathbf{a}_c/\|\mathbf{a}_c\|$.
The collective thrust is obtained by projecting the commanded acceleration onto
the \emph{current} body $z$-axis,
\begin{equation}
f_{\text{coll}} = m\,\mathbf{a}_c^\top (R\mathbf{e}_3),
\label{eq:thrust-proj}
\end{equation}
which prevents acceleration in the wrong direction under large attitude errors.

\section{INDI Inner Loop}\label{sec:indi}
The inner loop enforces the commanded angular acceleration without an aerodynamic model.
Over a short interval $\Delta t$,
\begin{equation}
\Delta\mathbf{y} = \mathbf{y}_{k+1}-\mathbf{y}_k
               = B\,(\mathbf{u}_{k+1}-\mathbf{u}_k)
               + (\mathbf{d}_{k+1}-\mathbf{d}_k),
\end{equation}
where $\mathbf{y}$ is the measured angular acceleration proxy and $B$ the
local control-effectiveness matrix.
Since aerodynamic disturbances $\mathbf{d}$ change slowly compared to the control rate,
$\Delta\mathbf{d}\!\approx\!0$ and
\[
\Delta\mathbf{y} \approx B\,\Delta\mathbf{u}.
\]
Thus the control update
$\mathbf{u}_{k+1}=\mathbf{u}_k+B^{-1}(\mathbf{y}_{\text{des}}-\mathbf{y}_k)$
implicitly cancels the aerodynamic contribution without modeling it.

\subsection{Why Aerodynamic Terms Are Irrelevant}
Let the true dynamics be $\dot{\mathbf{y}}=f(\mathbf{y},\mathbf{u})+\mathbf{d}$.
Classical NDI must model $\mathbf{d}$ explicitly.
INDI instead computes the input increment
\[
\Delta\mathbf{u} = B^{-1}\big(\dot{\mathbf{y}}_{\text{des}}-\dot{\mathbf{y}}_k\big),
\]
using measured acceleration differences.
Because $\mathbf{d}$ is approximately constant over $\Delta t$,
its effect cancels:
\[
(\mathbf{d}_{k+1}-\mathbf{d}_k)\approx 0.
\]
Consequently, all quasi-steady aerodynamic forces and moments—lift, induced drag,
or wing interference—are suppressed automatically.
The controller reacts only to *changes* in measured acceleration, making an explicit
aerodynamic model unnecessary while retaining high robustness.

\section{Controller Structure and Relation to Prior Work}\label{sec:controller-summary}
The resulting control stack can be interpreted as a two-level INDI cascade:
an outer geometric controller with incremental feedforward of aerodynamic
disturbances and an inner INDI loop realizing angular dynamics.
This configuration matches the most accurate setup in
Sun~et al.~\cite{Sun2021}, where coupling a flatness-based outer loop with
an INDI inner loop reduced position-tracking error by roughly 78 \% at
flight speeds up to 20 m/s compared to pure differential-flatness control.

\section{Summary}\label{sec:ctrl-summary}
The architecture preserves the intuitive structure of a geometric quadrotor
controller while inheriting the disturbance-rejection and aerodynamic robustness
of INDI.  Quasi-steady aerodynamic effects are naturally eliminated through
incremental updates, requiring no explicit aerodynamic model.

% --- References used in this chapter ---
% \cite{Sieberling2010,vanKampen2018,Tzoumanikas2021,Goodarzi2015,Mellinger2011,Sun2021}
